\frontheading{Abstract}
Prediction-market prices invite a probability interpretation, but a discrepancy between prices is not sufficient evidence of an executable arbitrage. Contract definitions, quoted liquidity, fees, observation timing and incomplete execution determine which apparent opportunities survive. This study develops a read-only research engine for Kalshi that combines explicit payoff-state checks, fixed-point order-book reconstruction, quantity-dependent acquisition costs and chronological probability evaluation.

The historical sample comprises 270 daily maximum-temperature events and 1,620 contracts across New York, Chicago and Miami between 19 July and 16 October 2026. Public REST collection provides 56,006 hourly candles. The primary noon specification retains 268 complete six-contract event panels: 162 for development, 53 for validation and 53 for later evaluation. A separate live pilot records 720 individual books across 36 contracts in 20 batch responses over approximately three minutes. Synthetic experiments test known payoff inconsistencies and the consequences of incomplete multi-leg execution.

There are 34 positive gross historical basket-side gaps among 536 conditional checks, but none remains positive under the stated current-fee scenario. The historical records do not establish executable depth or fills. In the live pilot, 378 basket checks have sufficient displayed depth and no positive net gap; 342 fail the depth requirement. Forty-six complete two-sided event observations support statewise optimisation, with no positive rechecked candidate. A conservative full-precision payoff audit and exceptional settlement provisions further limit unconditional arbitrage certification.

A one-parameter probability calibration fits a development temperature of 1.0344 but performs worse than the unchanged market-derived baseline on validation. The selection rule therefore retains the baseline, whose later noon log loss is 0.718589 nats and multiclass Brier score is 0.408273. Synthetic execution stresses show that positive expected value can coexist with losses from incomplete baskets. The contribution is a reproducible framework for distinguishing probability disagreement, conditional payoff gaps and executable opportunities, rather than evidence of realised trading profitability. The discussion records acquisition failures, contract-audit corrections, model rejection and remaining data-access and execution limitations.

\vspace{4pt}
\textbf{Keywords:} prediction markets; market microstructure; probability calibration; arbitrage screening; order-book liquidity; execution risk.
